% ============================================================
% ENGLISH TENSES — PAST PERFECT
% Version graphique LaTeX/TikZ
% Compilation : XeLaTeX
% ============================================================
\documentclass[a4paper,10pt]{article}

\usepackage[a4paper,top=9mm,bottom=8mm,left=11mm,right=11mm]{geometry}
\usepackage{fontspec}
\usepackage[french]{babel}
\usepackage{microtype}
\usepackage{tikz}
\usetikzlibrary{arrows.meta,calc,positioning}
\usepackage[most]{tcolorbox}
\usepackage{tabularx}
\usepackage{array}
\usepackage{xcolor}
\usepackage[normalem]{ulem}

\setmainfont{Noto Sans}
\setsansfont{Noto Sans}

% ------------------------------------------------------------
% PALETTE
% ------------------------------------------------------------
\definecolor{Green}{HTML}{356FA8}
\definecolor{GreenDark}{HTML}{244D78}
\definecolor{GreenLight}{HTML}{EAF2F9}
\definecolor{GreenLine}{HTML}{9BBBD7}

\definecolor{Blue}{HTML}{356FA8}
\definecolor{BlueLight}{HTML}{EAF2F9}

\definecolor{Red}{HTML}{D52F35}
\definecolor{RedLight}{HTML}{FCEBED}

\definecolor{Gold}{HTML}{B27A00}
\definecolor{GoldLight}{HTML}{FFF3D7}

\definecolor{Ink}{HTML}{172A32}
\definecolor{Grey}{HTML}{64757D}
\definecolor{Border}{HTML}{B9D3E3}
\definecolor{White}{HTML}{FFFFFF}

\pagestyle{empty}
\setlength{\parindent}{0pt}
\setlength{\parskip}{0pt}
\setlength{\tabcolsep}{3pt}

% ------------------------------------------------------------
% OUTER FRAME
% Left corners square, right corners rounded.
% The green stripe is independent and perfectly vertical.
% ------------------------------------------------------------
\newcommand{\PageFrame}{%
\begin{tikzpicture}[remember picture,overlay]
  \coordinate (A) at ([xshift=8mm,yshift=-8mm]current page.north west);
  \coordinate (B) at ([xshift=-8mm,yshift=-8mm]current page.north east);
  \coordinate (C) at ([xshift=-8mm,yshift=8mm]current page.south east);
  \coordinate (D) at ([xshift=8mm,yshift=8mm]current page.south west);
  \def\r{2.5mm}

  % Fine border with rounded right corners only
  \draw[Border,line width=.65pt]
    (A)
    -- ([xshift=-\r]B)
    arc[start angle=90,end angle=0,radius=\r]
    -- ([yshift=\r]C)
    arc[start angle=0,end angle=-90,radius=\r]
    -- (D) -- cycle;

  % Green vertical liseré
  \draw[Green,line width=3.2mm] (A) -- (D);
\end{tikzpicture}%
}

% ------------------------------------------------------------
% BOX STYLES
% ------------------------------------------------------------
\tcbset{
  enhanced,
  boxrule=.6pt,
  arc=2.8mm,
  left=4mm,right=4mm,top=2mm,bottom=2mm,
  coltext=Ink
}

\newtcolorbox{timelinebox}{
  colback=White,colframe=GreenLine
}

\newtcolorbox{greenbox}{
  colback=GreenLight,colframe=GreenLine
}

\newtcolorbox{bluebox}{
  colback=BlueLight,colframe=Border
}

\newtcolorbox{redbox}{
  colback=RedLight,colframe=Border
}

\newtcolorbox{whitebox}{
  colback=White,colframe=Border
}

\newtcolorbox{goldbox}{
  colback=GoldLight,colframe=GoldLight,
  boxrule=0pt,arc=2.5mm,
  left=2.7mm,right=2.7mm,top=1.25mm,bottom=1.25mm
}

% ------------------------------------------------------------
% SMALL LABELS
% ------------------------------------------------------------
\newcommand{\pill}[2][GreenLight]{%
  \tikz[baseline=(X.base)]{
    \node[
      fill=#1,
      rounded corners=2pt,
      inner xsep=4.2pt,
      inner ysep=1.6pt,
      font=\small\bfseries,
      text=Ink
    ] (X) {#2};
  }%
}

\newcommand{\sectiontitle}[2]{%
  {\large\bfseries\textcolor{#1}{#2}}%
}

\newcommand{\signalpill}[1]{%
  \tikz[baseline=(X.base)]{
    \node[
      fill=GoldLight,
      rounded corners=2pt,
      inner xsep=3.7pt,
      inner ysep=1.1pt,
      font=\scriptsize\bfseries,
      text=Ink
    ] (X) {#1};
  }%
}

% ------------------------------------------------------------
% TIMELINE
% ------------------------------------------------------------
\newcommand{\Timeline}{%
\begin{tikzpicture}[x=1cm,y=1cm]
  % main line
  \draw[-{Latex[length=3mm]},very thick,Ink] (0,0) -- (15.3,0);

  % NOW arrow
  \draw[-{Latex[length=3mm,width=2mm]},very thick,Green]
    (9.3,1.10) -- (9.3,.10);
  \node[font=\bfseries\small,text=Green] at (9.3,1.32) {NOW};

  % present marker
  \fill[Green] (9.3,0) circle (3.2pt);

  % finished action in the past
  \draw[Green,line width=4.5pt] (2.8,0) -- (4.65,0);
  \fill[Green] (4.65,0) circle (2.2pt);

  % labels
  \node[font=\scriptsize\bfseries,text=Grey] at (2.2,-.48) {PAST};
  \node[font=\scriptsize\bfseries,text=Green] at (9.3,-.48) {PRESENT};
  \node[font=\scriptsize\bfseries,text=Grey] at (13.3,-.48) {FUTURE};

  % action label
  \node[
    fill=GreenLight,rounded corners=2pt,
    inner xsep=5pt,inner ysep=2pt,
    font=\scriptsize\bfseries,text=GreenDark
  ] at (4.2,.72) {action finished before another past moment};
\end{tikzpicture}%
}

\begin{document}
\PageFrame

% ============================================================
% HEADER
% ============================================================
\vspace*{1mm}

\begin{center}
  {\fontsize{25}{27}\selectfont\bfseries\textcolor{Ink}{PAST PERFECT}}

  \vspace{0.75mm}
  {\large\textcolor{Grey}{Le past perfect anglais}}

  \vspace{3mm}

  \begin{tcolorbox}[
    width=.88\textwidth,
    colback=GreenLight,
    colframe=GreenLight,
    boxrule=0pt,
    arc=3mm,
    halign=center,
    top=2.2mm,bottom=2.2mm
  ]
    {\large\bfseries\textcolor{GreenDark}{SUBJECT + HAD + PAST PARTICIPLE}}
    \quad
    {\small(+ \textbf{-s / -es} à la 3\ieme{} personne du singulier)}
  \end{tcolorbox}
\end{center}

\vspace{1mm}

% ============================================================
% TIMELINE
% ============================================================
\begin{timelinebox}
  \sectiontitle{Green}{La ligne du temps}

  \vspace{0.75mm}
  \begin{center}
    \Timeline
  \end{center}

  \vspace{-1mm}
  \begin{center}
    \small
    Le \textbf{Past Perfect} décrit une action
    \textbf{\textcolor{Green}{qui s'était produite avant un autre moment du passé}}.
  \end{center}
\end{timelinebox}

\vspace{1.35mm}

% ============================================================
% AFFIRMATIVE / NEGATIVE
% ============================================================
\begin{minipage}[t]{.485\textwidth}
\begin{greenbox}
  \sectiontitle{Green}{Affirmative}

  \vspace{0.75mm}
  \begin{tabular}{@{}>{\bfseries}l l@{}}
    I / You / He / She / It / We / They & had worked.\\
    
  \end{tabular}

  \vspace{0.75mm}
  \textit{I had finished before 8 pm.}\\[-1mm]
  \textit{They had already left.}
\end{greenbox}
\end{minipage}
\hfill
\begin{minipage}[t]{.485\textwidth}
\begin{redbox}
  \sectiontitle{Red}{Negative}

  \vspace{0.75mm}
  \begin{tabular}{@{}>{\bfseries}l l@{}}
    I / You / He / She / It / We / They & \textcolor{Red}{\bfseries had not} worked.\\
    
  \end{tabular}

  \vspace{0.75mm}
  \textit{I hadn't finished.}\\[-1mm]
  \textit{They hadn't arrived.}
\end{redbox}
\end{minipage}

\vspace{1.35mm}

% ============================================================
% QUESTION
% ============================================================
\begin{bluebox}
  \sectiontitle{Blue}{Question}

  \vspace{0.75mm}
  \begin{tabularx}{\textwidth}{@{}>{\color{Blue}\bfseries}l X@{}}
    Had & I / you / he / she / it / we / they worked ?\\
    
  \end{tabularx}

  \vspace{0.75mm}
  \textit{Had you finished before 8 pm? \qquad Had she left?}

  \vspace{1mm}
  \begin{tcolorbox}[
    colback=RedLight,colframe=RedLight,boxrule=0pt,
    arc=2mm,left=2.7mm,right=2.7mm,top=1.25mm,bottom=1.25mm
  ]
    \textcolor{Red}{\bfseries !\quad Attention :}
    après \textbf{had}, on utilise le \textbf{past participle} :
    \textbf{She had worked.} \quad et non \quad
    \textcolor{Grey}{\sout{She had work.}}
  \end{tcolorbox}
\end{bluebox}

\vspace{1.35mm}

% ============================================================
% USES
% ============================================================
\begin{greenbox}
  \sectiontitle{Green}{When do we use the Simple Past?}

  \vspace{0.75mm}

  \begin{tabularx}{\textwidth}{@{}p{.31\textwidth} X@{}}
    \pill{EARLIER PAST ACTION}
      &
      Action qui s'est produite avant un autre moment passé.
      \newline \textit{I had eaten before I left.}\\[0.7mm]

    \pill{CAUSE / RESULT IN THE PAST}
      &
      Une action passée explique une situation ou un résultat ultérieur.
      \newline \textit{She was tired because she had worked by the time.}\\[0.7mm]

    \pill{SEQUENCE OF PAST EVENTS}
      &
      Le Past Perfect marque l'événement le plus ancien.
      \newline \textit{When I arrived, they had already left.}\\[0.7mm]

    \pill{BEFORE A PAST TIME}
      &
      Action terminée avant un moment précis du passé.
      \newline \textit{By 8 pm, he had finished the work.}
  \end{tabularx}
\end{greenbox}

\vspace{1.35mm}

% ============================================================
% EXAMPLES
% ============================================================
\begin{bluebox}
  \sectiontitle{Blue}{Examples}

  \vspace{0.75mm}
  \begin{tabularx}{\textwidth}{@{}X@{\hspace{7mm}}X@{}}
    \textbf{I had finished my work before dinner.}
    &
    \textbf{Paul had already gone home.}\\[-.5mm]
    \textit{J'ai visité Londres l'été dernier.}
    &
    \textit{Paul a travaillé tard hier.}\\[0.7mm]
    \textbf{She had left after I arrived.}
    &
    \textbf{They had eaten before the film started.}\\[-.5mm]
    \textit{Elle a ouvert la porte et s'est assise.}
    &
    \textit{As-tu regardé le film ?}
  \end{tabularx}
\end{bluebox}

\vspace{1.35mm}

% ============================================================
% SIGNAL WORDS
% ============================================================
\begin{goldbox}
  \begin{center}
    {\normalsize\bfseries\textcolor{Gold}{SIGNAL WORDS}}
    \hspace{3mm}
    \signalpill{already}\hspace{1mm}
    \signalpill{just}\hspace{1mm}
    \signalpill{before}\hspace{1mm}
    \signalpill{after}\hspace{1mm}
    \signalpill{by the time}\hspace{1mm}
    \signalpill{until then}

    \vspace{0.75mm}

    \signalpill{after}\hspace{1mm}
    \signalpill{recently}
  \end{center}
\end{goldbox}

\end{document}
